\usepackage{geometry}
\geometry{
    paperwidth=140mm,
    paperheight=203mm,
    left=2cm,
    right=2cm,
    top=2cm,
    bottom=1cm
    }

\usepackage{multirow}
\usepackage{amsmath}
\usepackage{ulem}
\usepackage{enumitem}
\usepackage{fancyhdr}
\pagestyle{fancy}
\renewcommand{\headrulewidth}{0pt}  % 添加这行，将横线宽度设为0
\fancyhf{}
\cfoot{-~\thepage~-} 
\setlength{\footskip}{14pt} 

% 自定义一个类似词典条目的环境
\newenvironment{entry}
{\begin{description}[
    labelwidth=2cm,       % 标签占据的宽度，可根据最长标签调整
    leftmargin=2cm,       % 整个条目左侧缩进（labelwidth + labelsep）
    labelsep=0.2cm,       % 标签和正文之间的间距
    align=right,           % 标签左对齐
    labelindent=0pt,      % 标签无额外缩进
    font=\normalfont,     % 标签字体正常（不加粗）
    itemindent=0pt,       % 正文首行不缩进
    ]}
{\end{description}}

\usepackage{tikz}
% 定义四角标命令 \cornernum{汉字}{左上}{右上}{左下}{右下}
\newcommand{\cornernum}[5]{%
    \begin{tikzpicture}[baseline=(word.base)]
        % 节点显示汉字
        \large
        \node[inner sep=0pt, outer sep=0pt] (word) {#1};
        % 标注四个角的数字，使用 tiny 字体或更小
        \node[anchor=south west, inner sep=0.1pt, font=\scriptsize] at (word.north west) {#2};
        \node[anchor=south east, inner sep=0.1pt, font=\scriptsize] at (word.north east) {#3};
        \node[anchor=north west, inner sep=0.1pt, font=\scriptsize] at (word.south west) {#4};
        \node[anchor=north east, inner sep=0.1pt, font=\scriptsize] at (word.south east) {#5};
    \end{tikzpicture}%
}